% =====================================================================
% Identidade visual PEE/COPPE/UFRJ
% Derivado do template de G. Gleizer (tema Dresden + logos institucionais),
% adaptado para 16:9 e com posicionamento dos logos via TikZ.
% =====================================================================

\mode<presentation>{
  \usetheme{Dresden}
  \usecolortheme{seagull}
}

% O cabecalho de navegacao do Dresden e removido; a orientacao fica por
% conta do slide de agenda e dos separadores de secao.
\setbeamertemplate{headline}{}
\setbeamertemplate{navigation symbols}{}
\setbeamertemplate{frametitle}[default][center]

% ---------------------------------------------------------------------
% Paleta
% ---------------------------------------------------------------------
\definecolor{peeblue}{RGB}{1,123,165}
\definecolor{peedark}{RGB}{0,74,99}
\definecolor{peegray}{RGB}{95,105,110}
\definecolor{peeaccent}{RGB}{200,90,40}

\setbeamercolor*{palette primary}{fg=black,bg=peeblue!80}
\setbeamercolor*{palette secondary}{fg=black,bg=peeblue!80!gray!80}
\setbeamercolor*{palette tertiary}{fg=black,bg=peeblue!100}
\setbeamercolor*{palette quaternary}{fg=black,bg=peeblue!110}

\setbeamercolor{frametitle}{fg=peedark}
\setbeamercolor{title}{fg=peedark}
\setbeamercolor{structure}{fg=peeblue}
\setbeamercolor{alerted text}{fg=peeaccent}

% ---------------------------------------------------------------------
% Rodape
%
% O rodape de duas linhas do Dresden monta caixas cuja soma excede
% \paperwidth em 16:9, produzindo um Overfull \hbox por slide. Aqui ele e
% substituido por uma barra unica cujas larguras somam exatamente 1.
% ---------------------------------------------------------------------
\setbeamertemplate{footline}{%
  \leavevmode%
  \hbox{%
    \begin{beamercolorbox}[wd=.34\paperwidth,ht=2.4ex,dp=1.2ex,left]{author in head/foot}%
      \hspace*{1.2ex}\usebeamerfont{author in head/foot}\insertshortauthor
    \end{beamercolorbox}%
    \begin{beamercolorbox}[wd=.46\paperwidth,ht=2.4ex,dp=1.2ex,center]{title in head/foot}%
      \usebeamerfont{title in head/foot}\insertshorttitle
    \end{beamercolorbox}%
    \begin{beamercolorbox}[wd=.20\paperwidth,ht=2.4ex,dp=1.2ex,right]{date in head/foot}%
      \usebeamerfont{date in head/foot}\insertshortinstitute
      \hspace*{1.2em}\insertframenumber{}\hspace*{1.2ex}
    \end{beamercolorbox}}%
  \vskip0pt%
}

% ---------------------------------------------------------------------
% Tipografia e elementos de lista
% ---------------------------------------------------------------------
\setbeamerfont{frametitle}{size=\large,series=\bfseries}
\setbeamerfont{title}{size=\Large,series=\bfseries}
\setbeamertemplate{itemize item}{\color{peeblue}$\blacktriangleright$}
\setbeamertemplate{itemize subitem}{\color{peegray}\textendash}
\setbeamertemplate{enumerate item}{\color{peeblue}\insertenumlabel.}
\setbeamercolor{caption name}{fg=peedark}
\setbeamerfont{caption}{size=\scriptsize}

% Blocos com cabecalho institucional
\setbeamertemplate{blocks}[rounded][shadow=false]
\setbeamercolor{block title}{fg=white,bg=peeblue}
\setbeamercolor{block body}{bg=peeblue!8}
\setbeamercolor{block title alerted}{fg=white,bg=peeaccent}
\setbeamercolor{block body alerted}{bg=peeaccent!8}

% ---------------------------------------------------------------------
% Logos de fundo (TikZ: posicionamento absoluto, imune ao aspect ratio)
% ---------------------------------------------------------------------
\newcommand{\peelogomarca}{%
  \begin{tikzpicture}[remember picture,overlay]
    \node[anchor=north west,inner sep=0pt]
      at ([shift={(0.28cm,-0.22cm)}]current page.north west)
      {\includegraphics[height=0.085\paperheight,keepaspectratio]{imgs/pee-logo-short.png}};
  \end{tikzpicture}%
}

\newcommand{\peelogoscapa}{%
  \begin{tikzpicture}[remember picture,overlay]
    \node[anchor=north west,inner sep=0pt]
      at ([shift={(0.9cm,-0.7cm)}]current page.north west)
      {\includegraphics[height=0.135\paperheight,keepaspectratio]{imgs/pee-logo.png}};
    \node[anchor=north east,inner sep=0pt]
      at ([shift={(-0.9cm,-0.7cm)}]current page.north east)
      {\includegraphics[height=0.135\paperheight,keepaspectratio]{imgs/coppe-logo.pdf}};
  \end{tikzpicture}%
}

% A marca dagua discreta acompanha todos os slides de conteudo.
\usebackgroundtemplate{\peelogomarca}

% A area util recua nas duas margens: o titulo centralizado fica livre da
% marca dagua no canto superior esquerdo sem estourar a caixa.
\setbeamersize{text margin left=1.1em,text margin right=1.1em}

% ---------------------------------------------------------------------
% Separador de secao automatico
% ---------------------------------------------------------------------
\AtBeginSection[]{%
  \begin{frame}[plain,noframenumbering]
    \vfill
    \centering
    \begin{beamercolorbox}[sep=10pt,center,shadow=false,rounded=true]{title}
      \usebeamerfont{title}\color{peedark}\insertsectionhead\par
    \end{beamercolorbox}
    \vfill
  \end{frame}%
}
